\documentclass[12pt,a4paper]{article}

% Paquetes básicos
\usepackage[utf8]{inputenc}
\usepackage[T1]{fontenc}
\usepackage[spanish]{babel}
\usepackage{amsmath, amssymb}
\usepackage{siunitx} % Ángulos
\usepackage{graphicx}
\usepackage{float}
\usepackage{geometry}
\geometry{margin=2.5cm}
\usepackage{hyperref}
\usepackage{caption}
\usepackage{subcaption}
\usepackage{setspace}
\usepackage{multirow}
\usepackage{minted} % Códigos
\graphicspath{{Figuras/}}
\onehalfspacing


% Datos
\author{Nombre y Apellido \\ Legajo: XXXXX}
\date{\today}

\begin{document}

% Carátula personalizada
\begin{titlepage}
    \centering
    
\includegraphics[width=0.25\textwidth]{escudo.png} 
    \hspace{1cm}
    
\includegraphics[width=0.3\textwidth]{download.png} \\[1cm]

    {\large \textbf{UNIVERSIDAD NACIONAL DE CÓRDOBA} \\[0.5cm]}
    {\large \textbf {FACULTAD DE CIENCIAS EXACTAS, FÍSICAS Y NATURALES} \\[1cm]}
    {\Large ELECTRÓNICA INDUSTRIAL \\[1cm]}

    \rule{\textwidth}{0.5pt} \\[0.5cm]
    {\Large \textbf{Trabajo Práctico Teórico N°4} \\[0.3cm]}
    {\Large \textbf{MOSFETs de Potencia} \\[0.3cm]}
    \rule{\textwidth}{0.5pt} \\[1cm]

\begin{tabular}{l l}
\textbf{Profesor Titular:} & Ing. Agüero Adrián.\\
\textbf{Profesor Adjunto:} & Ing. Gangi O. Sergio.\\
\textbf{Profesor Adjunto:} & Ing. Laboret Sergio.\\

\\
\textbf{Alumno:} Monja, Ernesto Joaquín. \\
\textbf{DNI:} 43.873.728 \\
\end{tabular}\\[3cm]

{\centering Año 2026 \par}

\end{titlepage}



% Índice
\tableofcontents
\newpage
\indent



\section{Introducción}
En este trabajo práctico se busca investigar sobre los MOSFETS de potencia existentes en el mercado, analizar sus hojas de datos y realizar simulaciones sobre los mismos en cuanto a parámetros característicos, gráficas y selección de los mismos para aplicaciones de la Electrónica de Potencia.


\section{Desarrollo}
\subsection{Lectura de Hojas de Datos}
\subsubsection{Aplicación al MOSFET IRF840}
El transistor de potencia elegido para este apartado es el IRF840 encapsulado SOT78 (TO220AB), el cual ya fue utilizado en el \href{https://github.com/ErnestMonja/Electronica-Industrial/blob/main/TPT1%20%20%20-%20Introducci%C3%B3n%20a%20la%20Electr%C3%B3nica%20de%20Potencia/Monja_TPT1_EI2026.pdf}{Trabajo Práctico N°1} de esta asignatura. Si revisamos su hoja de datos en internet, podemos encontrar las siguientes tablas:

\begin{figure}[H]
    \centering
    \includegraphics[width=0.9\linewidth]{Figura 1.png}
    \caption{\href{https://www.alldatasheet.com/datasheet-pdf/view/17804/PHILIPS/IRF840.html}{Tabla 2 y 3 de la hoja de datos del IRF840.}}
    \label{fig:fig_1}
\end{figure}
\begin{figure}[H]
    \centering
    \includegraphics[width=0.9\linewidth]{Figura 2.png}
    \caption{\href{https://www.alldatasheet.com/datasheet-pdf/view/17804/PHILIPS/IRF840.html}{Tabla 1 de la hoja de datos del IRF840.}}
    \label{fig:fig_2}
\end{figure}

De aquí se deducen los siguientes valores:
\begin{itemize}
    \item \underline{Corriente de Drenador ($I_{D}$):} El fabricante provee 3 valores máximos:
    \begin{table}[H]
        \centering
        \begin{tabular}{|c|c|c|} \hline
            Nombre & Valor de Corriente & Condiciones \\ \hline
            Corriente Continua de Drenador & $I_{D} = 8.5 \ [A]$ & $T_{amb} = 25 \ [°C]$; $V_{GS} = 10 \ [V]$  \\ \hline
            Corriente Continua de Drenador & $I_{D} = 5.4 \ [A]$ & $T_{amb} = 100 \ [°C]$; $V_{GS} = 10 \ [V]$  \\ \hline
            Corriente Pulsante de Drenador & $I_{DM} = 34 \ [A]$ & $T_{amb} = 25 \ [°C]$  \\ \hline
        \end{tabular}
        \caption{Valores de corriente de drenador del \href{https://www.alldatasheet.com/datasheet-pdf/view/17804/PHILIPS/IRF840.html}{IRF840}.}
        \label{tab:tab_1}
    \end{table}
    Sin embargo, nos quedaremos con el valor de: $I_{D} = 8.5 \ [A]$.

    \item \underline{Tensión Drenador-Surtidor ($V_{DS}$):} Se corresponde con el valor de tensión de ruptura de la unión que forman el surtidor y el drenador. En este caso es igual a: $V_{DSS} = 500 \ [V]$ dentro de un rango de $T_{j} = 25-100 \ [°C]$.

    \item \underline{Resistencia Drenador-Surtidor de Conducción ($r_{DS}$):} Es la resistencia que ofrece el MOSFET cuando está totalmente encendido. Aquí es igual a: $R_{DS(ON)} = 0.6 \ [\Omega]$ para una $V_{GS} = 10 \ [V]$ e $I_{D} = 4.8 \ [A]$.

    \item \underline{Tensión de Compuerta-Surtidor Límite ($V_{GS(TO)}$):} Es el voltaje mínimo que se debe aplicar entre la compuerta y el surtidor para que el MOSFET empiece a conducir. En este caso es igual a: $V_{GS(TO)} = 3 \ [V]$ con $V_{DS} = V_{GS}$ e $I_{D} = 0.25 \ [mA]$.

     \item \underline{Tensión de Compuerta-Surtidor Límite ($V_{GS}$):} Indica el rango de tensión máximo (positivo y negativo) que puede soportar la compuerta y se tiene que esta es igual a: $V_{GS} = \pm 30 \ [V]$.

     \item \underline{Corriente de Perdida de la Compuerta ($I_{G}$):} Es una corriente de fuga muy pequeña que atraviesa la compuerta, la cual es igual a: $I_{G} = 10 \ [nA]$ con $V_{GS} = \pm 30 \ [V]$ y $V_{DS} = 0 \ [V]$.

     \item \underline{Potencia Disipada ($P_{D}$):} Se trata de la potencia máxima que el componente puede disipar en forma de calor antes de quemarse y es igual a: $P_{D} = 147 \ [W]$ con una $T_{mb} = 25 \ [°C]$.
\end{itemize}

Si exploramos más, se encuentra tambien las gráficas de Área de Operación Segura (o Safe Operation Area del acrónimo del inglés) la cual se presenta a continuación:
\begin{figure}[H]
    \centering
    \includegraphics[width=0.6\linewidth]{Figura 3.png}
    \caption{\href{https://www.alldatasheet.com/datasheet-pdf/view/17804/PHILIPS/IRF840.html}{SOA del IRF840.}}
    \label{fig:fig_3}
\end{figure}

Finalmente se identifican las curvas de capacidades parásitas de entrada ($C_{iss}$), salida ($C_{oss}$) y de realimentación ($C_{rss}$) las cuales se presentan a continuación:
\begin{figure}[H]
    \centering
    \includegraphics[width=0.5\linewidth]{Figura 4.png}
    \caption{\href{https://www.alldatasheet.com/datasheet-pdf/view/17804/PHILIPS/IRF840.html}{Capacidades parásitas del IRF840.}}
    \label{fig:fig_4}
\end{figure}

Aquí se tienen que:
\begin{itemize}
    \item \underline{Capacidad de Entrada ($C_{iss} = C_{GS} + C_{GD}$):} Es la más alta de todas en valor nominal y determina la rapidez con la que puede encenderse el MOSFET, ya que antes de tener una tensión $V_{GS}$, se debe cargar este capacitor.
    \item \underline{Capacidad de Salida ($C_{oss} = C_{DS} + C_{GD}$):} Se trata de la capacidad de salida del transistor y resulta importante para calcular las pérdidas en la conmutación.
    \item \underline{Capacidad de Transferencia Inversa ($C_{rss} = C_{GD}$):} También conocida como la capacitancia de Miller, es la más crítica ya que retroalimenta la salida con la entrada durante la conmutación, frenando así la conmutación.
\end{itemize}

Se observa que las capacidades parásitas disminuyen al aumentar la tensión drenador-surtidor ($V_{DS}$). Esto se debe al ensanchamiento de las zonas de deplexión internas del semiconductor. La capacidad de entrada ($C_{iss}$) es la más estable, mientras que la de salida ($C_{oss}$) y transferencia ($C_{rss}$) muestran una dependencia no lineal muy marcada, reduciendo su valor significativamente a voltajes altos, lo que favorece menores tiempos de conmutación en regímenes de alta tensión.



\subsubsection{Conclusión Parcial sobre la Lectura de Datos}
En esta consigna, se ha podido identificar satisfactoriamente de la hoja de datos:
\begin{itemize}
    \item Los parámetros eléctricos típicos de interés de un MOSFET de Potencia  y su encapsulado.
    \item Se ha identificado la gráfica de Área de Operación Segura.
    \item Se ha identificado la gráfica de las capacidades parásitas del MOSFET, se ha explicado su significado y también el motivo de su variación.
\end{itemize}

Por lo tanto se da como concluida, esta primera parte del este Trabajo Práctico.



\subsection{Simulación de Conmutación}
\subsubsection{Aplicación para el IRF150}
En esta sección, se propone realizar una simulación en LTSpice del siguiente circuito: 
\begin{figure}[H]
    \centering
    \includegraphics[width=0.63\linewidth]{Figura 5.png}
    \caption{Simulación en LTSpice de la conmutación de un IRF150.}
    \label{fig:fig_5}
\end{figure}

Dado que no se contaba con el modelo del IRF150 en LTSpice, se opto por tomar el modelo de LTSpice propuesto por \href{https://github.com/drahnr/oregano/blob/master/data/models/IRF150.model}{Bernhard Schuster} (Véase el hipervínculo dentro del nombre). De esta simulación, se pueden obtener las siguientes gráficas:
\begin{figure}[H]
    \centering
    \includegraphics[width=0.8\linewidth]{Figura 6.png}
    \caption{Simulación en LTSpice de la conmutación de un IRF150.}
    \label{fig:fig_6}
\end{figure}

Aquí se identifican:
\begin{itemize}
    \item \underline{Curva de Tensión Compuerta-Surtidor ($V_{GS}$):} Como la curva verde.
    \item \underline{Curva de la Corriente del Drenador ($I_{D}$):} Como la curva azul.
    \item \underline{Curva de Potencia en la Conmutación ($P_{D}$):} Como la curva roja.
\end{itemize}

Si graficamos la curva de $V_{GS}$ (en verde) y la curva de $V_{DS}$ (en azul), se podrán identificar los tiempos característicos de la conmutación del MOSFET:
\begin{itemize}
    \item \underline{Tiempo de Retraso de Encendido ($t_{d(on)}$) y Tiempo de Subida ($t_{r}$):} Se presenta a continuación 2 imágenes compiladas en una, donde: a la izquierda se tiene una medición del $t_{d(on)} \approx 117.92 \ [ns]$ y a la derecha se tiene una medición del $t_{r} \approx 445.03 \ [ns]$ tal como se ve en la siguiente gráfica:
    \begin{figure}[H]
        \centering
        \includegraphics[width=0.9\linewidth]{Figura 7.png}
        \caption{Simulación en LTSpice de la conmutación de un IRF150.}
        \label{fig:fig_7}
    \end{figure}
    
    \item \underline{Tiempo de Retraso de Apagado ($t_{d(off)}$) y Tiempo de Caída ($t_{f}$):} Se presenta a continuación 2 imágenes compiladas en una, donde: a la izquierda se tiene una medición del $t_{d(off)} \approx 197.84 \ [ns]$ y a la derecha se tiene una medición del $t_{f} \approx 695.68 \ [ns]$ tal como se ve en la siguiente gráfica:
    \begin{figure}[H]
        \centering
        \includegraphics[width=0.9\linewidth]{Figura 8.png}
        \caption{Simulación en LTSpice de la conmutación de un IRF150.}
        \label{fig:fig_8}
    \end{figure}
\end{itemize}



\subsubsection{Conclusión Parcial de la Simulación}
Las primeras simulaciones de tensión, corriente y potencia son las gráficas a esperar que se estudiaron en el teórico. Se observa el proceso de carga del capacitor de entrada en la Figura (\ref{fig:fig_6}) en la gráfica de la tensión $V_{GS}$ (Véase la curva verde). Además, se pudo observar la existencia de un pico de potencia durante la conmutación así mismo como las pérdidas en conducción del transistor.

En cuanto a la medición de los tiempos característicos del MOSFET, se tiene que estos no se corresponden a los que provee la hoja de datos:
\begin{figure}[H]
    \centering
    \includegraphics[width=0.7\linewidth]{Figura 9.png}
    \caption{\href{https://www.alldatasheet.com/datasheet-pdf/view/37003/SAMSUNG/IRF150.html}{Tabla 2 de la hoja de datos del IRF150.}}
    \label{fig:fig_9}
\end{figure}

Esto se debe principalmente a que el fabricante provee un entorno de trabajo distinto al utilizado para obtener estas mediciones, es decir que el fabricante midió estos tiempos con distintos valores de tensiones, corrientes, etc. los cuales indudablemente afectarán a los resultados de las simulaciones realizadas. No tan solo eso, sino que se debe considerar que el modelo utilizado en LTSpice no siempre es exacto al modelo real del MOSFET. Estas consideraciones explican la diferencia entre valores simulados y reales del IRF150.

\section{Bibliografía}

\begin{itemize}
    \item \href{https://www.alldatasheet.com/datasheet-pdf/view/17804/PHILIPS/IRF840.html}{Hoja de datos del IRF840.}
    \item \href{https://github.com/drahnr/oregano/blob/master/data/models/IRF150.model}{Repositorio de Bernhard Schuster con el Modelo de LTSpice del IRF150.}
    \item \href{https://www.alldatasheet.com/datasheet-pdf/view/37003/SAMSUNG/IRF150.html}{Hoja de datos del IRF150.}
\end{itemize}

\end{document}
